\documentclass[11pt,a4paper]{article}

\usepackage[utf8]{inputenc}
\usepackage[T1]{fontenc}
\usepackage[provide=*,english]{babel}
\usepackage{lmodern}
\usepackage{microtype}
\usepackage{amsmath,amssymb,mathtools}
\usepackage{geometry}
\usepackage{xcolor}
\usepackage{enumitem}
\usepackage{fancyhdr}
\usepackage[hidelinks]{hyperref}

\geometry{left=2.3cm,right=2.3cm,top=2.2cm,bottom=2.25cm,headheight=14pt}
\setlength{\parindent}{0pt}
\setlength{\parskip}{0.42em}
\setlist[itemize]{topsep=0.25em,itemsep=0.18em,leftmargin=1.5em}
\allowdisplaybreaks

\definecolor{blue}{RGB}{34,73,110}
\definecolor{gray}{RGB}{90,96,102}

\pagestyle{fancy}
\fancyhf{}
\fancyhead[L]{\small\color{gray}GC2D: ABBA4 implicit}
\fancyhead[R]{\small\color{gray}Jacobian formula summary}
\fancyfoot[C]{\thepage}
\renewcommand{\headrulewidth}{0.3pt}

\newcommand{\Id}{I}
\newcommand{\R}{\mathbb{R}}
\newcommand{\norm}[1]{\left\lVert#1\right\rVert}

\begin{document}

\begin{center}
  {\LARGE\bfseries\color{blue}ABBA4 Implicit Jacobian Formulas}\par
  \vspace{0.35em}
  {\large Local implicit-function factors and the triple-jump product}\par
  \vspace{0.45em}
  {\small Reduced physical branch for non-autonomous GC2D dynamics}
\end{center}

\section{Composition notation}

Let
\[
 \gamma=\frac{1}{2-2^{1/3}},
 \qquad
 \delta=-\frac{2^{1/3}}{2-2^{1/3}},
 \qquad
 (c_1,c_2,c_3)=(\gamma,\delta,\gamma).
\]
For outer step size $h$, define
\[
 \ell_j=c_jh,
 \qquad
 \tau_j=t_n+h\sum_{r<j}c_r,
 \qquad
 s_j=\frac{\ell_j}{2},
 \qquad
 j=1,2,3.
\]
The signed physical substeps are
\[
 z^{[0]}=z_n,
 \qquad
 z^{[j]}=\Psi^{[2]}_{\ell_j,\tau_j}(z^{[j-1]}),
 \qquad
 z_{n+1}=z^{[3]}.
\]
Every quantity below with first index $j$ belongs to one independent
projection solve.  In particular, the three multipliers $\mu_j$ and their
stage points must not be shared.

\section{Stage Jacobians for one signed substep}

Let
\[
 W_{j,i}=D_zf(\text{stage point }i,\text{stage time }i),
 \qquad i=1,2,3,4,
\]
and define
\[
 S_j=W_{j,2}+W_{j,3}.
\]
The matrix-product order must be preserved.  The exact duplicated ABBA
Jacobian has $2\times2$ blocks
\[
 J_{\mathrm{ABBA},j}
 =\begin{pmatrix}A_j&B_j\\C_j&D_j\end{pmatrix},
\]
where
\begin{align*}
 A_j&=\Id_2+s_j^2W_{j,4}S_j,\\
 B_j&=s_j(W_{j,1}+W_{j,4})
      +s_j^3W_{j,4}S_jW_{j,1},\\
 C_j&=s_jS_j,\\
 D_j&=\Id_2+s_j^2S_jW_{j,1}.
\end{align*}
These expressions are valid without modification for the negative middle
substep because $s_2=\delta h/2<0$.

\section{Residual derivatives \texorpdfstring{$K_j$}{Kj} and \texorpdfstring{$L_j$}{Lj}}

For the reduced projection residual
\[
 R_j(z,\mu)
 =G\Phi^{\mathrm{ABBA}}_{\ell_j,\tau_j}(Ez+G^T\mu)+2\mu,
\]
define
\[
 K_j=D_\mu R_j,
 \qquad
 L_j=D_zR_j.
\]
Their compact block expressions are
\[
 \boxed{
 K_j=A_j-B_j-C_j+D_j+2\Id_2,
 }
\]
\[
 \boxed{
 L_j=A_j+B_j-C_j-D_j.
 }
\]
Equivalently,
\[
 \boxed{
 \begin{aligned}
 K_j={}&4\Id_2-s_j(W_{j,1}+W_{j,2}+W_{j,3}+W_{j,4})\\
 &+s_j^2\bigl(W_{j,4}S_j+S_jW_{j,1}\bigr)
 -s_j^3W_{j,4}S_jW_{j,1},
 \end{aligned}
 }
\]
and
\[
 \boxed{
 \begin{aligned}
 L_j={}&s_j(W_{j,1}-W_{j,2}-W_{j,3}+W_{j,4})\\
 &+s_j^2\bigl(W_{j,4}S_j-S_jW_{j,1}\bigr)
 +s_j^3W_{j,4}S_jW_{j,1}.
 \end{aligned}
 }
\]

The converged multiplier is a function
\[
 \mu_j^*=\mu_{\ell_j,\tau_j}(z^{[j-1]}).
\]
Differentiating
\(
 R_j(z,\mu_j^*(z))=0
\)
gives
\[
 L_j+K_jM_j=0,
 \qquad
 M_j:=D_z\mu_j^*,
\]
and hence
\[
 \boxed{M_j=-K_j^{-1}L_j.}
\]
An implementation solves $K_jM_j=-L_j$ and does not form the inverse.

\section{Physical Jacobian of one substep}

Define
\[
 P_j=A_j+B_j,
 \qquad
 Q_j=A_j-B_j+\Id_2.
\]
In expanded form,
\[
 P_j=\Id_2+s_j(W_{j,1}+W_{j,4})
 +s_j^2W_{j,4}S_j+s_j^3W_{j,4}S_jW_{j,1},
\]
\[
 Q_j=2\Id_2-s_j(W_{j,1}+W_{j,4})
 +s_j^2W_{j,4}S_j-s_j^3W_{j,4}S_jW_{j,1}.
\]
The exact ideal-root physical Jacobian is
\[
 \boxed{
 J_j
 :=D_z\Psi^{[2]}_{\ell_j,\tau_j}(z^{[j-1]})
 =P_j-Q_jK_j^{-1}L_j.
 }
\]
To construct it stably, solve the two-right-hand-side system
\[
 K_jT_j=L_j,
 \qquad
 J_j=P_j-Q_jT_j.
\]
The Newton matrix $K_j$ is not the physical Jacobian $J_j$.

\section{Equivalent stage-increment factorization}

For one substep, the physical mean is
\[
 z^{[j]}=z^{[j-1]}
 +\frac{\ell_j}{4}(k_{j,1}+k_{j,2}+k_{j,3}+k_{j,4}).
\]
The total derivatives, including $M_j=D\mu_j^*$, obey
\begin{align*}
 Dk_{j,1}&=W_{j,1}(\Id_2-M_j),\\
 Du_j^{(1)}&=\Id_2+M_j+s_jDk_{j,1},\\
 Dk_{j,2}&=W_{j,2}Du_j^{(1)},\\
 Dk_{j,3}&=W_{j,3}Du_j^{(1)},\\
 Dv_j^{(2)}&=\Id_2-M_j+s_j(Dk_{j,2}+Dk_{j,3}),\\
 Dk_{j,4}&=W_{j,4}Dv_j^{(2)}.
\end{align*}
Therefore,
\[
 \boxed{
 J_j=\Id_2+\frac{\ell_j}{4}
 \sum_{i=1}^4Dk_{j,i}.
 }
\]
This stage-increment recurrence and the implicit-function formula are two
factorizations of the same local tangent.  Omitting $M_j$ differentiates a
different, explicitly averaged map.

\section{Complete ABBA4 Jacobian}

The outer chain rule gives
\[
 \boxed{
 D_z\Psi^{[4]}_{h,t_n}(z_n)=J_3J_2J_1.
 }
\]
The chronological order matters: $J_1$ acts first and $J_3$ multiplies from
the left.  No derivative of the Newton or Broyden iteration history appears;
the formulas differentiate the converged implicit roots.

For a tangent vector $\xi_n$, apply each local Jacobian without constructing
the full outer matrix.  Set $\eta_0=\xi_n$ and, for $j=1,2,3$, compute
\[
 b_j=L_j\eta_{j-1},
 \qquad
 K_jy_j=b_j,
 \qquad
 \boxed{\eta_j=P_j\eta_{j-1}-Q_jy_j.}
\]
Then
\[
 \xi_{n+1}=\eta_3.
\]
Each application adds one $2\times2$ linear solve per particle and no new
nonlinear solve.

\section{Accumulation over many outer steps}

Let $\mathcal J_n$ be the local ABBA4 Jacobian for outer step $n$, and let
$\mathcal M_n$ be the derivative of the numerical state with respect to the
initial condition.  Initialize
\[
 \mathcal M_0=\Id_2
\]
and update
\[
 \boxed{
 \mathcal M_{n+1}=\mathcal J_n\mathcal M_n,
 \qquad
 \mathcal J_n=J_{n,3}J_{n,2}J_{n,1}.
 }
\]
Thus the most recent local Jacobian always multiplies from the left.

For the ideal projected roots,
\[
 J_{n,j}^TJ_0J_{n,j}=J_0
\]
for every signed substep, and consequently
\[
 \mathcal J_n^TJ_0\mathcal J_n=J_0,
 \qquad
 \mathcal M_n^TJ_0\mathcal M_n=J_0.
\]
In two dimensions these identities are equivalent to unit determinant.  A
finite nonlinear tolerance makes the emitted numerical map differ slightly
from the ideal-root map evaluated by these formulas.

\section{Implementation checklist}

For each accepted outer step:
\begin{itemize}
 \item retain the ordered three substep snapshots and their independent
       multipliers;
 \item evaluate $W_{j,1},\ldots,W_{j,4}$ at the signed stage times;
 \item form $K_j,L_j,P_j,Q_j$ without commuting matrix products;
 \item solve $K_jT_j=L_j$ or apply the tangent through one right-hand side;
 \item compose in the order $J_3J_2J_1$;
 \item compare against centered finite differences only at a scale larger than
       the three nonlinear residual floors.
\end{itemize}

\end{document}
